% Revisão consolidada: 25/11/2026. Não antecipa resultados finais.
\chapter{Estado da arte}
\section{Método e delimitação do levantamento}
O levantamento desta versão é narrativo e complementar, conduzido sob responsabilidade de Caroline. Foram examinados primeiro os arquivos enviados e os documentos do acervo do DoseCerta; depois, consultaram-se fontes públicas acadêmicas e oficiais para completar os fundamentos e a comparação. Não se apresenta esta busca como revisão sistemática ou como seleção exaustiva da literatura.
Os eixos utilizados são: forma de organizar e liberar doses, interação e acessibilidade, rastreabilidade, integração com o cuidador, verificação do evento e limites de evidência. Documentos do projeto descrevem decisões do DoseCerta; não são evidência independente de superioridade frente a outras soluções.
\pendencia{EA02-REV}{Caroline: concluir a revisão bibliográfica, registrar bases, datas, expressões de busca, critérios e referências adicionais. Se for adotado um método sistemático, executar e documentar o protocolo antes de assim denominar o levantamento.}
\section{Organizadores e dispensadores}
Organizadores passivos, dispensadores programáveis e sistemas conectados diferem no controle sobre a disponibilização da dose e no tipo de registro gerado. Essas categorias são usadas como eixos conceituais, e não como uma classificação comercial exaustiva. Comparar apenas presença de alarmes não esclarece reabastecimento, dose unitária, tratamento de falha ou qualidade da evidência de uso.
A revisão de Salmensuu e colaboradores avaliou dispensadores para pessoas idosas que vivem em casa em domínios de experiência, saúde, uso de serviços, custos e equidade. Foram reunidos 25 estudos originais a partir de estudos e revisões incluídos. Os autores apontam evidência moderada para alguns resultados de saúde, mas não evidência forte e uniforme dos benefícios em todos os domínios. \cite{SALMENSUU2025}
Assim, os achados sustentam investigar sistemas de apoio, mas não permitem atribuir ao DoseCerta os efeitos de equipamentos ou intervenções diferentes. Tampouco justificam prometer redução de custos sem avaliação própria. \cite{SALMENSUU2025}
\section{Verificação do evento e monitoramento}
O registro de abertura de compartimento ou de comando de motor tem significado diferente de confirmação de liberação, retirada ou ingestão. A proposta DoseCerta separa leitura da cavidade, execução do mecanismo e inferências da visão. A contribuição pretendida está na integração rastreável, cuja avaliação depende de evidências específicas de cada etapa. \cite{DC_COMPILADO}
\pendencia{EA03-COMP}{Completar o quadro comparativo com modelos específicos e fontes primárias: capacidade, autonomia, dimensões, preço com data/região, reabastecimento e mecanismo. Não preencher parâmetros comerciais ausentes por analogia nem afirmar vantagem sem ensaio comparável.}
\section{Interação digital e agente}
O Care organiza o acompanhamento e as configurações; o agente pretende tornar a consulta e certas solicitações acessíveis por linguagem natural. A arquitetura documentada separa uma interface conversacional de voz, sem ações MCP, e um agente de teste com oito ferramentas e aprovação de escritas. Essa separação delimita o estado documental da integração, em vez de descrever uma assistência autônoma já demonstrada. \cite{DC_APP,DC_FOUNDRY}
Acessibilidade deve ser investigada em tarefas concretas, inclusive compreensão de falhas e confirmações. Adotar orientações das WCAG 2.2 como referência não substitui avaliação com o público pretendido nem autoriza afirmar conformidade integral. \cite{WCAG22}
\section{Posicionamento do DoseCerta}
O trabalho propõe articular módulos de medicamento removíveis, separação unitária, controle local, rastreabilidade digital, observação visual e assistência conversacional. A demonstração de utilidade dessa integração permanece distinta da demonstração individual de componentes. Nenhuma superioridade, eficácia clínica ou economia é concluída nesta etapa.
\pendencia{EA04-COMP}{Inserir comparação empírica somente após definição de população, cenário, desfechos e dados. Deixar explícitos os parâmetros ainda não medidos e as diferenças entre as soluções comparadas.}
